\documentclass[10pt,a4paper]{amsart}
\usepackage[margin=19mm]{geometry}
\usepackage{amsmath,amssymb,mathtools}
\usepackage[scr=rsfs,cal=euler]{mathalfa}
\usepackage{xfrac}
\usepackage[unicode,hidelinks,bookmarks=false]{hyperref}
\newcommand{\ie}{i.e.\ }
\newcommand{\cf}{cf.\ }
\newcommand{\resp}{respectively\ }
\newcommand{\Subsection}{\subsection}
\providecommand{\nobreakdash}{\nobreak\hbox{-}}
\setlength{\emergencystretch}{2em}
\multlinegap0pt
\usepackage{bm}
\usepackage{bbm}
\usepackage{dsfont}
\usepackage{xspace}
\usepackage{cancel}
\usepackage{mathtools}
\usepackage{amsthm}
\makeatletter
\@ifundefined{theorem}{\newtheorem{theorem}{Theorem}}{}
\@ifundefined{lemma}{\newtheorem{lemma}[theorem]{Lemma}}{}
\@ifundefined{remark}{\newtheorem{remark}[theorem]{Remark}}{}
\makeatother
\newcommand{\N}{\mathbb{N}}
\newcommand{\Per}{\text{Per}}
\title[Synchronizing Dynamical Systems: Finitely presented systems ]{Synchronizing Dynamical Systems: Finitely presented systems and Ruelle algebras}
\author{Robin J. Deeley, Andrew M. Stocker}
\date{Source: arXiv:2501.13909v2 (CC BY 4.0)}
\begin{document}
\maketitle

\noindent\emph{Source and licence.} Synchronizing Dynamical Systems: Finitely presented systems and Ruelle algebras. Version arXiv:2501.13909v2 (CC BY 4.0), distributed under the Creative Commons Attribution 4.0 International licence, as recorded in the arXiv licence metadata for this exact version and shown on the abstract page https://arxiv.org/abs/2501.13909v2. Full rights notice: licenses/arxiv-2501.13909-CC-BY-4.0.txt.\par\smallskip

\noindent\textit{Editorial scope.} Counted unit: the Remark whose source label is \texttt{None} in the article (source0.tex lines 697--699, source label \texttt{None}), delivered together with the article's own complete original proof of it (source0.tex lines 700--731, one \texttt{proof} environment of 32 lines). No mathematics is written, repaired or re-derived: statement and proof are the article's own text line for line. Required context retained uncounted: lem:periodic-h-implies-equals (lines 209--211). Every definition, display and label of those lines is retained unchanged. Omitted from this package and not redistributed: 1--208, 212--696, 732--1237. Nothing inside a retained excerpt is omitted or altered: every retained body is delivered exactly as the article's lines are written, so no omission receipt of any kind accompanies this package. Citations: the retained excerpts carry no citation command at all, so no bibliography entry of the article is transcribed and no reference is redistributed. Reference adaptation: none was needed and none was made; every reference inside the delivered unit resolves inside it, so no unresolved reference or citation is produced. Printed slips kept literal: no printed word, symbol, hypothesis or number of the counted statement or of its proof was altered. Layout and class: the article's own class is \texttt{amsart} with packages including amsfonts, amsmath, amsthm, amssymb, amscd, latexsym, enumerate, etoolbox, mathrsfs, comment, xy, thmtools, thm-restate, tikz-cd; the wrapper is the lane's pinned \texttt{amsart} editorial wrapper at 10pt on A4, which restates the theorem-like environments the retained text needs and adds the article's own macro definitions that the retained text needs (\texttt{\textbackslash N}, \texttt{\textbackslash Per}), with extra wrapper packages (bm, bbm, dsfont, xspace, cancel, mathtools). The delivered numbering is the unit's own. Assets: no retained span contains a figure environment, a graphics-inclusion command, a PDF-page inclusion, a table or a TikZ picture; no image file is copied into this package, the package declares no asset dependency and no image file is redistributed with it. Licence: version arXiv:2501.13909v2 of this article is distributed under the Creative Commons Attribution 4.0 International licence, as recorded in the arXiv licence metadata for this exact version and shown on the abstract page https://arxiv.org/abs/2501.13909v2. A full rights notice is delivered at licenses/arxiv-2501.13909-CC-BY-4.0.txt. Characters: every retained excerpt is pure ASCII, so the delivered engine, XeLaTeX with its default Latin Modern text font, reports no missing character.
\begin{lemma}\label{lem:periodic-h-implies-equals}
If $(X,\varphi)$ is an expansive dynamical system and $p , q \in \Per(X,\varphi)$ are periodic points such that either $p \sim_\text{s} q$ or $p \sim_\text{u} q$, then $p = q$.  In particular $p \sim_\text{h} q$ implies $p = q$.
\end{lemma}

\begin{remark}
It is worth noting that this result is stronger than the statement that $\pi|_{X^u(Q)}$ is one-to-one; the element $y$ has unique preimage with respect to the map $\pi : X \rightarrow Y$.
\end{remark}

\begin{proof}
Let $y\in Y^u(P)$. Then, by Lemma \ref{lem:periodic-h-implies-equals},  $y \in Y^u(p)$ for a unique $p\in P$. By assumption, there is unique $q\in X$ such that $\pi^{-1}(p)=\{ q \}$. Moreover, $q$ is a periodic point since $p$ is. Let $M>1$ be such that 
\[
\varphi^M(q)=q \hbox{ and }f^M(p)=p.
\]

Since $\pi$ is onto, $\pi^{-1}(y)$ is non-empty. Take $x\in \pi^{-1}(y)$ and form the sequence in $X$, $(\varphi^{M\cdot n}(x))_{n\in \N}$. Let $L$ be a limit point of this sequence. Then $\pi(L)=\lim_{k\rightarrow \infty} \pi( \varphi^{M\cdot n_k}(x))$ for some subsequence $(\varphi^{M\cdot n_k}(x))_{k\in \N}$. Using the definition of a factor map, $y\in Y^u(p)$, and $f^M(p)=p$, we get
\[
\pi(L)=\lim_{k\rightarrow \infty} \pi( \varphi^{M\cdot n_k}(x))=\lim_{k\rightarrow \infty} f^{M\cdot n_k}(\pi(x))=\lim_{k\rightarrow \infty}f^{M\cdot n_k}(y)=p.
\]
Hence $L=q$ because $\pi^{-1}(p)=\{q\}$. Using the fact that $X$ is compact and the fact that all limit points of the sequence are $q$, we have that $\lim_{n \rightarrow \infty} \varphi^{M\cdot n}(x)=q$. 

Next, we show $x \in X^u(q)$. Let $\epsilon>0$. Since $X$ is compact and $\varphi$ is continuous, $\varphi$ is uniformly continuous. Hence there exists $\delta>0$ such that 
\[
d(\varphi^r(a), \varphi^r(b))<\epsilon
\]
whenever $d(a, b)<\delta$ and $r=0, 1, \ldots, M-1$.

Since $\lim_{n \rightarrow \infty} \varphi^{M\cdot n}(x)=q$, there exists $N\in \N$ such that 
\[
d(\varphi^{M\cdot n}(x), q) < \delta \hbox{ for }n\ge N.
\]
Take $k \geq M \cdot N$. We have that
\[
d(\varphi^k(x), \varphi^k(q))=d(\varphi^{M\cdot l + r}(x), \varphi^{M\cdot l + r}(q))=d(\varphi^r(\varphi^{M\cdot l}(x)), \varphi^r(q))
\]
for some $l\ge N$ and $0\le r \le M-1$. Using our uniformly continuous condition above, we have that 
\[ d(\varphi^r(\varphi^{M\cdot l}(x)), \varphi^r(q))< \epsilon \]
because $d((\varphi^{M\cdot l}(x), q)< \delta$. Hence $x\in X^u(q)$.

Finally, since $\pi$ is u-resolving and $\emptyset \neq \pi^{-1}(y) \subseteq X^u(q)$, we have that $\pi^{-1}(y)$ is a single point in $X^u(q) \subseteq X^u(Q)$.
\end{proof}

\end{document}
